\documentclass[12pt]{article}
\usepackage{../TopicChapters/ceai}


\begin{document}


\title{CE 488/5588 Homework 01\\
\vspace{2in}
\pmb{Name: \rule{2in}{0.15mm}}
\author{}
}
\maketitle

\doublespacing
\newcommand*{\I}{\imath} % imaginary unit i

\newpage
\section*{P1. Read and Reflection}
Read this report \href{https://reports.weforum.org/docs/WEF_Future_of_Jobs_Report_2025.pdf}{Future of Jobs Report} from WEF, and comment the following questions:

\begin{itemize}
    \item How would the skill set shift or reskilling need happen by 2030? 
    \item Would we have net engineering job creation by 2030, based on this report due to AI and automation
\end{itemize}

\section*{P2. Multiple Choice and Explain}

\subsection*{Question 1}
Which process below is \textbf{NOT} a deductive reasoning process?

\begin{enumerate}
    \item A climate scientist uses historical records of the last 100 years and infers a conclusion that warming is unavoidable in the 21st century.
    \item A climate scientist models the atmosphere using physical laws based on partial differential equations.
    \item A student uses Newton's laws and calculates how long it takes when dropping a ball from 10 ft above the ground.
    \item A hydraulic engineer models river flooding based on fluid mechanics using physics-based simulation software.
\end{enumerate}
Explain your answer:\\
\fullunderline\\ 
\fullunderline \\ 
\fullunderline \\
\fullunderline




\subsection*{Question 2}
Which process below is \textbf{NOT} inductive?

\begin{enumerate}
    \item Tornado has not hit inner Kansas City. Therefore, it would never hit Kansas City.
    \item A student does a data-fitting problem and uses a linear model.
    \item A transportation engineer monitors the traffic for a crossing. Using AM and PM traffic data collected over many days, a model is built to predict AM and PM traffic levels.
    \item An environmental engineer applies a classical pollution equation by plugging air monitoring data to estimate pollution levels.
\end{enumerate}

Explain your answer:\\
\fullunderline\\ 
\fullunderline \\ 
\fullunderline \\
\fullunderline

\newpage
\section*{P3. Programming Exercises}

\paragraph{General Guide:} In this class, although we have a slogan - \textit{no coding but vibe programming}, we will accommodate different needs. In all programming practices, code snippets will be provided: you have the right to do the hardcore approach, coding without using Gemini, or the most modern free-style approach, totally completing the AI programming by prompting and interacting with Gemini. 

\subsection*{Exercise 1: Creating a Time Series Plot}
Write a Python program to create a time series using the $\sin$ function and plot the curve. The $x$-axis should be labeled \textbf{seconds}, and the $y$-axis should be labeled \textbf{acceleration}.

\textbf{Report and print out your plot. }

\textit{Hints:}
\begin{itemize}
    \item Use \texttt{numpy} for generating the time series data with \texttt{np.sin()}.
    \item Use \texttt{matplotlib.pyplot} to plot the curve.
\end{itemize}

\textbf{Example Code Outline:}
\begin{verbatim}
import numpy as np
import matplotlib.pyplot as plt

# Generate time series data
time = np.linspace(0, 10, 100)  # 0 to 10 seconds, 100 points
acceleration = np.sin(time)

# Plotting
plt.plot(time, acceleration)
plt.xlabel('Seconds')
plt.ylabel('Acceleration')
plt.title('Time Series of Acceleration')
plt.show()
\end{verbatim}

\subsection*{Exercise 2: Generating and Plotting a Histogram}
Write a Python program to generate a random Gaussian dataset with 100 data points using \texttt{numpy.random} and plot its histogram.

\textbf{Report and print out your plot. }

\textit{Hints:}
\begin{itemize}
    \item Use \texttt{np.random.randn()} to generate Gaussian-distributed data.
    \item Use \texttt{plt.hist()} to plot the histogram.
\end{itemize}

\textbf{Example Code Outline:}
\begin{verbatim}
import numpy as np
import matplotlib.pyplot as plt

# Generate random data
data = np.random.randn(100)  # Gaussian distribution

# Plotting histogram
plt.hist(data, bins=10, edgecolor='black')
plt.xlabel('Value')
plt.ylabel('Frequency')
plt.title('Histogram of Random Gaussian Data')
plt.show()
\end{verbatim}

\end{document}
